\section{CSP: backtracking, revisi\'on hacia adelante, AC-3 y ordenamiento}

% TODO (20 pts implementación): describe brevemente tu abstracción de CSP
% (src/csp/problem.py) y cómo backtracking.py, forward_checking.py, ac3.py
% y heuristics.py encajan entre sí (¿AC3 solo como preprocesamiento o como
% MAC durante la búsqueda?).

\subsection{N-reinas}
% TODO (5+5 pts aplicar): N=8 y N=100 con backtracking. Función objetivo /
% formulación CSP (variables, dominios, restricciones). Tabla con nodos
% expandidos, tiempo, y si el ordenamiento (MRV/grado/LCV) cambió el resultado.
% Ver results/tables/nqueens.csv.

\subsection{Enumerar todas las soluciones, N=100}
% TODO (5 pts): reporta cuántas soluciones encontraste, en cuánto tiempo, y
% si la búsqueda fue exhaustiva o se cortó (ver
% results/tables/nqueens_enumeration.csv, campo extra_exhaustive). Discute
% qué tan grande es el espacio de soluciones reales para N=100 y por qué
% (o si) una enumeración exhaustiva es viable.

\subsection{Coloreado de grafos}
% TODO (5 pts aplicar): 50 y 1000 nodos con backtracking. Formulación CSP
% (variables=vértices, dominio=colores, restricción=adyacentes distintos) y
% cómo decidiste el valor de k. Ver results/tables/graph_coloring.csv.
